\documentclass[12pt,a4paper]{article}

% Packages
\usepackage{amsmath, amsfonts, amssymb, amsthm, bm, bbm, comment} 
\usepackage{enumerate, graphicx, xcolor, tcolorbox, hyphenat, circledsteps}
\usepackage[shortlabels]{enumitem}
\usepackage{url, hyperref}
\hypersetup{
	colorlinks   = true,
	urlcolor     = cyan,
	linkcolor    = red,
	citecolor    = black
}

% Formatting
\usepackage[a4paper, noheadfoot, scale={0.8,0.82}]{geometry}
\usepackage[T1]{fontenc}
\newcommand{\thetitle}{Clarity Pointers}
\newcommand{\theauthor}{Luca Bollini}
\newcommand{\thedate}{November 2026}
\hyphenpenalty=10000
\parindent0pt

% Commands
\newcommand{\N}{\mathbb{N}}
\newcommand{\Z}{\mathbb{Z}}
\newcommand{\Q}{\mathbb{Q}}
\newcommand{\R}{\mathbb{R}}
\newcommand{\eps}{\varepsilon}
\newcommand{\ddt}[1]{\frac{d}{d#1}}
\newcommand{\ddttwo}[1]{\frac{d^2}{d#1^2}}
\newcommand{\ts}{\thinspace}

% Colours
\newtcolorbox{questionbox}{
	colframe=black, 
	colback=white, 
	boxrule=0.2mm,
	arc=0mm,
}


% Specific packages
\usepackage{stmaryrd}
\usepackage{tikz}

\newcommand{\contradiction}{%
\begin{tikzpicture}[rotate=45,x=0.5ex,y=0.5ex]
\draw[line width=.2ex] (0,2) -- (3,2) (0,1) -- (3,1) (1,3) -- (1,0) (2,3) -- (2,0);
\end{tikzpicture}
}

\begin{document}

\framebox[16.8cm]{{\sc University of Warwick}\hfill{\sc Mathematics Institute}}

\vspace{2mm}
{\sc  \theauthor}\hfill{\sc \thedate}

\vspace{3mm}

\begin{center}
{\large \textbf {\underline{\thetitle}}}
\end{center}

\vspace{3mm}

\begin{enumerate}
    \item \underline{Write in full sentences.} They do not have to be complicated or eloquent. In fact, short, simple sentences are usually best!

    \item \underline{Use punctuation.} Your sentences should adhere to normal English grammar rules. That means: capital letter at the start, commas separating clauses, and a full stop at the end.    

    \item \underline{Avoid shorthand.} Especially $\therefore$ in place of ``therefore'', $\because$ in place of ``because'', and contradiction symbols such as $\lightning$ and $\contradiction$ can be difficult to decipher. 
    
    Blackboard shorthand such as $\exists$ for ``there exists'' and $\forall$ for ``for all'' are acceptable in assignments and exams, but are inappropriate in professional writing.

    \item \underline{Only use the ``implies'' symbol when you actually mean it!} The implies symbol is an ``if ... then ...'' statement. It should not be used to indicate the next line of a proof.

    \item \underline{Introduce new notation.} Do not assume your reader knows that $f$ is a function and $\eps$ is a positive real number. Write ``Let $f:X\to Y$ be a function'' and ``Let $\eps>0$''.
    
    \item \underline{Be precise.} When you see a statement such as
    \[
    \text{``The sequence $(a_n)$ satisfies } |a_n - l|<\eps.\text{''}
    \]
    you should be asking yourself: ``What values can $a_n$ take? For which $n$ is this true? For which $\eps$? Is $l$ a real number or an integer?'' You should specify that ``$(a_n)$ is a sequence of real numbers'' and the above statement is true ``for any $n>N$''.

    \item \underline{Check the order of mathematical clauses.} There is a stark difference between
    \[
    \text{``For any }N\in\Z,\text{ there exists }M\in\Z\text{ such that }M<N.\text{''}
    \]
    and
    \[
    \text{``There exists }M\in\Z\text{ such that }M<N, \text{ for any }N\in\Z.\text{''}
    \]
    One statement says you can always find a smaller integer, while the other says that there is a smallest integer!

    \item \underline{Proof read.} It is very easy to accidentally get the order of two statements the wrong way round, or misplace punctuation in a way that changes the meaning of a sentence. Check that what you have written actually communicates what you were intending.
\end{enumerate}
    







\end{document}